\documentclass[11pt]{article}
\usepackage[margin=1in]{geometry}
\usepackage{amsmath,amssymb,mathtools}
\usepackage{fontspec,unicode-math}
\setmainfont{Latin Modern Roman}
\setsansfont{Latin Modern Sans}
\setmathfont{Latin Modern Math}
\usepackage{xurl,hyperref}
\hypersetup{colorlinks=true,urlcolor=blue,linkcolor=blue}
\newcommand{\sref}[2]{\hyperlink{source:#1:#2}{[S#2]}}
\newcommand{\cataloguescope}[1]{\paragraph{Recorded mathematical question.}#1}
\begin{document}
\section{Distance exponent of random series-parallel graphs}\label{statement}

\subsection{Definitions and mathematical statement}

\paragraph{Shared notation.}$\mathbb N=\{1,2,\ldots\}$, and $\mathbb Z,\mathbb Q,\mathbb R,\mathbb C$ denote the integers, rationals, reals and complex numbers. Any other conventions are specified locally.


Starting with one source--sink edge $G_0$, replace every edge independently at each generation by two edges in series with probability $p$ or two parallel edges with probability $1-p$. Let $\Delta_n$ be the unit-edge source--sink distance in generation $n$. Equivalently,
\[\Delta_n\stackrel d=\begin{cases}\Delta_{n-1}^{(1)}+\Delta_{n-1}^{(2)},&\text{with probability }p,\\\min(\Delta_{n-1}^{(1)},\Delta_{n-1}^{(2)}),&\text{with probability }1-p,\end{cases}\]
with independent copies and an independent choice of case. Put $\delta(p)=\lim_{n\to\infty}n^{-1}\log\mathbb E\Delta_n$, the exponential mean-distance rate. The source proves this rate exists and is positive for $p>1/2$.
\paragraph{Statement recorded in the supplied source.}
Determine $p\mapsto\delta(p)$ on $[1/2,1]$, and decide whether $\delta(1/2)=0$. \sref{statement}{2}
\subsection{Short English statement}
Determine the exponential distance-growth rate in a random series-parallel graph, especially at its critical parameter.
\subsection{Sources}
\begin{description}
\item[\hypertarget{source:statement:1}{S1}] ULAM.AI and the UnsolvedMath Contributors, \emph{UnsolvedMath: A Curated Collection of Open Mathematics Problems}, catalogue entry. CC BY 4.0. \url{https://huggingface.co/datasets/ulamai/UnsolvedMath}.
\item[\hypertarget{source:statement:2}{S2}] Itai Benjamini, \emph{Euclidean vs. Graph Metric}, version dated 16 July 2012; original author manuscript preserved in Arquivo.pt capture of 31 December 2020, Section 9.4, final model and Question 9.6, p. 18. \url{https://arquivo.pt/noFrame/replay/20201231041538id_/http://www.wisdom.weizmann.ac.il/~itai/erd100.pdf}.
\end{description}
\label{end:statement}
\end{document}
